\section{Buildup finite element solution}
The next step is to solve this system and find the unknowns. This is done by using \emph{MATLAB Equation Solver}. For large No.
of elements we should take advantage of the sparse system property to save both memory and processor resources. This means that
we don't have to save the full stiffness matrix and \gls{loadVector}, but just nonzero diagonals in \gls{stiffnessMatrix} and load vector.
For the \gls{sparseMat} we also have access to much more efficient solver \glspl{algorithm}. For the problem in hand,
however, we have saved and solved the full system. \Cref{lst:Build_Sol} shows the MATLAB code for building up \gls{FEMSol} with
both linear and cubic \glspl{basisFunction}.
\begin{remark}
	The primary solution in some numerical methods\footnote{Such as \glspl{finiteDiffMethod}.} are the function values at
	predefined nodal points. In each case a smooth solution may then be constructed by interpolating these points. In FEM, on the
	other hand, interpolation is an integrated part of the method, see \cref{eqn:Truncated_FEM_Solution}.
\end{remark}	\text{}
For a discussion of these and other relevant materials consult \cite{Bertoti:1997,Carlson:1967,KarelRektorys:1975,Szeidl:2001}.